\chapter{Исследовательский раздел}

В этом разделе будет выполнено тестирование разработанного программного обеспечения на наборе PDF-файлов,
приложенных к лабораторной работе, а также проведено сравнение трёх реализаций, полученных с использованием
различных больших языковых моделей (ChatGPT / DeepSeek / Qwen).

Для проверки корректности работы использовался набор PDF-файлов:
\texttt{\_01.pdf ... \_10.pdf}, \texttt{main (2).pdf}, \texttt{ВКР Селез.pdf}.
Для каждой из трёх реализаций выполнялся запуск на одном и том же наборе входных файлов, после чего
формировался отдельный CSV-файл результатов.

\section{Форма представления результатов}

Для каждого файла фиксировались:
\begin{itemize}
  \item название PDF-файла;
  \item признак наличия ошибки (строго: \textbf{«найдена ошибка»} / \textbf{«ошибка не найдена»});
  \item координаты первого нарушения (страница, строка) либо «-», если координаты определить невозможно
  (например, если в документе не обнаружен раздел «СОДЕРЖАНИЕ/ОГЛАВЛЕНИЕ»).
\end{itemize}

\subsection{Результаты тестирования (реализация ChatGPT / GPT-5.2)}

\begin{longtable}{|p{0.32\textwidth}|p{0.30\textwidth}|p{0.24\textwidth}|}
\hline
\textbf{Название PDF} & \textbf{Признак} & \textbf{Координаты первого нарушения} \\
\hline
\endfirsthead
\hline
\textbf{Название PDF} & \textbf{Признак} & \textbf{Координаты первого нарушения} \\
\hline
\endhead

\texttt{\_01.pdf} & найдена ошибка & - \\ \hline
\texttt{\_02.pdf} & найдена ошибка & (2,1) \\ \hline
\texttt{\_03.pdf} & найдена ошибка & (2,1) \\ \hline
\texttt{\_04.pdf} & найдена ошибка & (5,1) \\ \hline
\texttt{\_05.pdf} & найдена ошибка & (1,2) \\ \hline
\texttt{\_06.pdf} & найдена ошибка & - \\ \hline
\texttt{\_07.pdf} & найдена ошибка & - \\ \hline
\texttt{\_08.pdf} & найдена ошибка & (1,1) \\ \hline
\texttt{\_09.pdf} & найдена ошибка & - \\ \hline
\texttt{\_10.pdf} & найдена ошибка & (2,1) \\ \hline
\texttt{main (2).pdf} & найдена ошибка & (5,1) \\ \hline
\texttt{ВКР Селез.pdf} & найдена ошибка & - \\ \hline

\end{longtable}

\subsection{Результаты тестирования (реализация DeepSeek DeepThink)}

\begin{longtable}{|p{0.32\textwidth}|p{0.30\textwidth}|p{0.24\textwidth}|}
\hline
\textbf{Название PDF} & \textbf{Признак} & \textbf{Координаты первого нарушения} \\
\hline
\endfirsthead
\hline
\textbf{Название PDF} & \textbf{Признак} & \textbf{Координаты первого нарушения} \\
\hline
\endhead

\texttt{\_01.pdf} & найдена ошибка & - \\ \hline
\texttt{\_02.pdf} & найдена ошибка & (2,1) \\ \hline
\texttt{\_03.pdf} & найдена ошибка & (2,1) \\ \hline
\texttt{\_04.pdf} & найдена ошибка & (5,1) \\ \hline
\texttt{\_05.pdf} & найдена ошибка & (1,2) \\ \hline
\texttt{\_06.pdf} & найдена ошибка & - \\ \hline
\texttt{\_07.pdf} & найдена ошибка & - \\ \hline
\texttt{\_08.pdf} & найдена ошибка & (1,1) \\ \hline
\texttt{\_09.pdf} & найдена ошибка & - \\ \hline
\texttt{\_10.pdf} & найдена ошибка & (2,1) \\ \hline
\texttt{main (2).pdf} & найдена ошибка & (5,1) \\ \hline
\texttt{ВКР Селез.pdf} & найдена ошибка & - \\ \hline

\end{longtable}

\subsection{Результаты тестирования (реализация Qwen Coder)}

\begin{longtable}{|p{0.32\textwidth}|p{0.30\textwidth}|p{0.24\textwidth}|}
\hline
\textbf{Название PDF} & \textbf{Признак} & \textbf{Координаты первого нарушения} \\
\hline
\endfirsthead
\hline
\textbf{Название PDF} & \textbf{Признак} & \textbf{Координаты первого нарушения} \\
\hline
\endhead

\texttt{\_01.pdf} & найдена ошибка & - \\ \hline
\texttt{\_02.pdf} & найдена ошибка & (2,1) \\ \hline
\texttt{\_03.pdf} & найдена ошибка & (2,1) \\ \hline
\texttt{\_04.pdf} & найдена ошибка & (5,1) \\ \hline
\texttt{\_05.pdf} & найдена ошибка & (1,2) \\ \hline
\texttt{\_06.pdf} & найдена ошибка & - \\ \hline
\texttt{\_07.pdf} & найдена ошибка & - \\ \hline
\texttt{\_08.pdf} & найдена ошибка & (1,1) \\ \hline
\texttt{\_09.pdf} & найдена ошибка & - \\ \hline
\texttt{\_10.pdf} & найдена ошибка & (2,1) \\ \hline
\texttt{main (2).pdf} & найдена ошибка & (5,1) \\ \hline
\texttt{ВКР Селез.pdf} & найдена ошибка & - \\ \hline

\end{longtable}

\section{Сравнение результатов трёх реализаций}

Для сравнения использовались три реализации, разработанные с применением различных БЯМ.
Результаты представлены в виде пары: \textit{статус, координаты}.

\begin{longtable}{|p{0.22\textwidth}|p{0.23\textwidth}|p{0.23\textwidth}|p{0.23\textwidth}|}
\hline
\textbf{Файл} & \textbf{ChatGPT} & \textbf{DeepSeek} & \textbf{Qwen} \\
\hline
\endfirsthead
\hline
\textbf{Файл} & \textbf{ChatGPT} & \textbf{DeepSeek} & \textbf{Qwen} \\
\hline
\endhead

\texttt{\_01.pdf} & найдена ошибка, - & найдена ошибка, - & найдена ошибка, - \\ \hline
\texttt{\_02.pdf} & найдена ошибка, (2,1) & найдена ошибка, (2,1) & найдена ошибка, (2,1) \\ \hline
\texttt{\_03.pdf} & найдена ошибка, (2,1) & найдена ошибка, (2,1) & найдена ошибка, (2,1) \\ \hline
\texttt{\_04.pdf} & найдена ошибка, (5,1) & найдена ошибка, (5,1) & найдена ошибка, (5,1) \\ \hline
\texttt{\_05.pdf} & найдена ошибка, (1,2) & найдена ошибка, (1,2) & найдена ошибка, (1,2) \\ \hline
\texttt{\_06.pdf} & найдена ошибка, - & найдена ошибка, - & найдена ошибка, - \\ \hline
\texttt{\_07.pdf} & найдена ошибка, - & найдена ошибка, - & найдена ошибка, - \\ \hline
\texttt{\_08.pdf} & найдена ошибка, (1,1) & найдена ошибка, (1,1) & найдена ошибка, (1,1) \\ \hline
\texttt{\_09.pdf} & найдена ошибка, - & найдена ошибка, - & найдена ошибка, - \\ \hline
\texttt{\_10.pdf} & найдена ошибка, (2,1) & найдена ошибка, (2,1) & найдена ошибка, (2,1) \\ \hline
\texttt{main (2).pdf} & найдена ошибка, (5,1) & найдена ошибка, (5,1) & найдена ошибка, (5,1) \\ \hline
\texttt{ВКР Селез.pdf} & найдена ошибка, - & найдена ошибка, - & найдена ошибка, - \\ \hline

\end{longtable}

\section{Оценка сложности}

Временная сложность алгоритма линейна по числу строк, извлечённых из PDF-документа,
то есть $O(N)$, где $N$ — количество строк после извлечения.
Дополнительная память также оценивается как $O(N)$ из-за хранения списка строк с координатами.

\section{Вывод}

По результатам тестирования все три реализации на данном наборе PDF-файлов сформировали одинаковые статусы
и совпадающие координаты первого нарушения в случаях, когда координаты можно определить.
Следовательно, различий в итоговом результате проверки для выбранного набора входных данных не выявлено.

В этом разделе было выполнено тестирование программ на приложенном наборе PDF-файлов,
построены таблицы результатов для каждой реализации и составлена сравнительная таблица,
а также дана оценка вычислительной сложности алгоритма.
